% ========================================================
% BÀI TẬP TOÁN 5 - KẾT NỐI TRI THỨC VỚI CUỘC SỐNG
% Bài 55: Luyện tập chung
% Định dạng: BTTeX Part III - Trắc nghiệm điền số ngắn
% ========================================================

% Câu 1
\begin{ex}
  Một khối kim loại hình lập phương có cạnh $4\text{ cm}$. Thể tích khối kim loại là bao nhiêu xăng-ti-mét khối?
  \shortans{64}
  \loigiai{
    Thể tích là 4 x 4 x 4 = 64 cm3.
  }
\end{ex}

% Câu 2
\begin{ex}
  Một khối kim loại hình lập phương có cạnh $5\text{ cm}$. Thể tích khối kim loại là bao nhiêu xăng-ti-mét khối?
  \shortans{125}
  \loigiai{
    Thể tích là 5 x 5 x 5 = 125 cm3.
  }
\end{ex}

% Câu 3
\begin{ex}
  Một khối kim loại hình lập phương có cạnh $6\text{ cm}$. Thể tích khối kim loại là bao nhiêu xăng-ti-mét khối?
  \shortans{216}
  \loigiai{
    Thể tích là 6 x 6 x 6 = 216 cm3.
  }
\end{ex}

% Câu 4
\begin{ex}
  Một khối kim loại hình lập phương có cạnh $7\text{ cm}$. Thể tích khối kim loại là bao nhiêu xăng-ti-mét khối?
  \shortans{343}
  \loigiai{
    Thể tích là 7 x 7 x 7 = 343 cm3.
  }
\end{ex}

% Câu 5
\begin{ex}
  Một khối kim loại hình lập phương có cạnh $8\text{ cm}$. Thể tích khối kim loại là bao nhiêu xăng-ti-mét khối?
  \shortans{512}
  \loigiai{
    Thể tích là 8 x 8 x 8 = 512 cm3.
  }
\end{ex}

% Câu 6
\begin{ex}
  Một khối kim loại hình lập phương có cạnh $9\text{ cm}$. Thể tích khối kim loại là bao nhiêu xăng-ti-mét khối?
  \shortans{729}
  \loigiai{
    Thể tích là 9 x 9 x 9 = 729 cm3.
  }
\end{ex}

% Câu 7
\begin{ex}
  Một khối kim loại hình lập phương có cạnh $10\text{ cm}$. Thể tích khối kim loại là bao nhiêu xăng-ti-mét khối?
  \shortans{1000}
  \loigiai{
    Thể tích là 10 x 10 x 10 = 1000 cm3.
  }
\end{ex}

% Câu 8
\begin{ex}
  Một khối kim loại hình lập phương có cạnh $11\text{ cm}$. Thể tích khối kim loại là bao nhiêu xăng-ti-mét khối?
  \shortans{1331}
  \loigiai{
    Thể tích là 11 x 11 x 11 = 1331 cm3.
  }
\end{ex}

% Câu 9
\begin{ex}
  Một khối kim loại hình lập phương có cạnh $12\text{ cm}$. Thể tích khối kim loại là bao nhiêu xăng-ti-mét khối?
  \shortans{1728}
  \loigiai{
    Thể tích là 12 x 12 x 12 = 1728 cm3.
  }
\end{ex}

% Câu 10
\begin{ex}
  Một khối kim loại hình lập phương có cạnh $13\text{ cm}$. Thể tích khối kim loại là bao nhiêu xăng-ti-mét khối?
  \shortans{2197}
  \loigiai{
    Thể tích là 13 x 13 x 13 = 2197 cm3.
  }
\end{ex}

% Câu 11
\begin{ex}
  Một khối kim loại hình lập phương có cạnh $14\text{ cm}$. Thể tích khối kim loại là bao nhiêu xăng-ti-mét khối?
  \shortans{2744}
  \loigiai{
    Thể tích là 14 x 14 x 14 = 2744 cm3.
  }
\end{ex}

% Câu 12
\begin{ex}
  Một khối kim loại hình lập phương có cạnh $15\text{ cm}$. Thể tích khối kim loại là bao nhiêu xăng-ti-mét khối?
  \shortans{3375}
  \loigiai{
    Thể tích là 15 x 15 x 15 = 3375 cm3.
  }
\end{ex}

% Câu 13
\begin{ex}
  Một khối kim loại hình lập phương có cạnh $16\text{ cm}$. Thể tích khối kim loại là bao nhiêu xăng-ti-mét khối?
  \shortans{4096}
  \loigiai{
    Thể tích là 16 x 16 x 16 = 4096 cm3.
  }
\end{ex}

% Câu 14
\begin{ex}
  Một khối kim loại hình lập phương có cạnh $17\text{ cm}$. Thể tích khối kim loại là bao nhiêu xăng-ti-mét khối?
  \shortans{4913}
  \loigiai{
    Thể tích là 17 x 17 x 17 = 4913 cm3.
  }
\end{ex}

% Câu 15
\begin{ex}
  Một khối kim loại hình lập phương có cạnh $18\text{ cm}$. Thể tích khối kim loại là bao nhiêu xăng-ti-mét khối?
  \shortans{5832}
  \loigiai{
    Thể tích là 18 x 18 x 18 = 5832 cm3.
  }
\end{ex}

% Câu 16
\begin{ex}
  Một khối kim loại hình lập phương có cạnh $19\text{ cm}$. Thể tích khối kim loại là bao nhiêu xăng-ti-mét khối?
  \shortans{6859}
  \loigiai{
    Thể tích là 19 x 19 x 19 = 6859 cm3.
  }
\end{ex}

% Câu 17
\begin{ex}
  Một khối kim loại hình lập phương có cạnh $20\text{ cm}$. Thể tích khối kim loại là bao nhiêu xăng-ti-mét khối?
  \shortans{8000}
  \loigiai{
    Thể tích là 20 x 20 x 20 = 8000 cm3.
  }
\end{ex}

% Câu 18
\begin{ex}
  Một khối kim loại hình lập phương có cạnh $21\text{ cm}$. Thể tích khối kim loại là bao nhiêu xăng-ti-mét khối?
  \shortans{9261}
  \loigiai{
    Thể tích là 21 x 21 x 21 = 9261 cm3.
  }
\end{ex}

% Câu 19
\begin{ex}
  Một khối kim loại hình lập phương có cạnh $22\text{ cm}$. Thể tích khối kim loại là bao nhiêu xăng-ti-mét khối?
  \shortans{10648}
  \loigiai{
    Thể tích là 22 x 22 x 22 = 10648 cm3.
  }
\end{ex}

% Câu 20
\begin{ex}
  Một khối kim loại hình lập phương có cạnh $23\text{ cm}$. Thể tích khối kim loại là bao nhiêu xăng-ti-mét khối?
  \shortans{12167}
  \loigiai{
    Thể tích là 23 x 23 x 23 = 12167 cm3.
  }
\end{ex}

% Câu 21
\begin{ex}
  Một khối kim loại hình lập phương có cạnh $24\text{ cm}$. Thể tích khối kim loại là bao nhiêu xăng-ti-mét khối?
  \shortans{13824}
  \loigiai{
    Thể tích là 24 x 24 x 24 = 13824 cm3.
  }
\end{ex}

% Câu 22
\begin{ex}
  Một khối kim loại hình lập phương có cạnh $25\text{ cm}$. Thể tích khối kim loại là bao nhiêu xăng-ti-mét khối?
  \shortans{15625}
  \loigiai{
    Thể tích là 25 x 25 x 25 = 15625 cm3.
  }
\end{ex}

% Câu 23
\begin{ex}
  Một khối kim loại hình lập phương có cạnh $26\text{ cm}$. Thể tích khối kim loại là bao nhiêu xăng-ti-mét khối?
  \shortans{17576}
  \loigiai{
    Thể tích là 26 x 26 x 26 = 17576 cm3.
  }
\end{ex}

% Câu 24
\begin{ex}
  Một khối kim loại hình lập phương có cạnh $27\text{ cm}$. Thể tích khối kim loại là bao nhiêu xăng-ti-mét khối?
  \shortans{19683}
  \loigiai{
    Thể tích là 27 x 27 x 27 = 19683 cm3.
  }
\end{ex}

% Câu 25
\begin{ex}
  Một khối kim loại hình lập phương có cạnh $28\text{ cm}$. Thể tích khối kim loại là bao nhiêu xăng-ti-mét khối?
  \shortans{21952}
  \loigiai{
    Thể tích là 28 x 28 x 28 = 21952 cm3.
  }
\end{ex}
